% Paper C: Bridging Infinite Domains with Hash-Consed Term IDs
% The "Bridge Paper" - How demand-driven term creation bridges finite and infinite
%
% For God so loved the world that he gave his only begotten Son,
% that whoever believes in him should not perish but have eternal life.
% John 3:16

\documentclass[11pt,a4paper]{article}

\usepackage[utf8]{inputenc}
\usepackage{amsmath,amssymb,amsthm}
\usepackage{algorithm}
\usepackage{algpseudocode}
\usepackage{booktabs}
\usepackage{graphicx}
\usepackage{hyperref}
\usepackage{listings}
\usepackage{xcolor}
\usepackage{xspace}
\usepackage{tikz}
\usetikzlibrary{shapes,arrows,positioning,fit,calc}

% Theorem environments
\newtheorem{theorem}{Theorem}
\newtheorem{lemma}[theorem]{Lemma}
\newtheorem{definition}[theorem]{Definition}
\newtheorem{proposition}[theorem]{Proposition}
\newtheorem{corollary}[theorem]{Corollary}

% Include shared definitions
% Shared Definitions for miniKanren 1-Bit Paper Suite
% Include with: % Shared Definitions for miniKanren 1-Bit Paper Suite
% Include with: % Shared Definitions for miniKanren 1-Bit Paper Suite
% Include with: \input{../shared_chirho/shared_defs_chirho}
%
% For God so loved the world that he gave his only begotten Son,
% that whoever believes in him should not perish but have eternal life.
% John 3:16

% ============================================================================
% Core Definitions (shared across all papers)
% ============================================================================

% The Finite-Domain miniKanren Kernel
\newcommand{\fdmk}{FD-miniKanren\xspace}  % Finite-Domain miniKanren

% ============================================================================
% Scope Clarification (include in each paper's introduction)
% ============================================================================

\newcommand{\scopeclarification}{%
\textbf{Scope Clarification.}
This paper addresses the \emph{Finite-Domain miniKanren Kernel} (\fdmk),
where variable domains are finite bitmasks.
Standard miniKanren unification over unbounded term structures uses a separate
reference implementation; the bit-parallel approach accelerates the constraint
propagation core, not full structural unification.
}

% ============================================================================
% Common Notation
% ============================================================================

% Domains
\newcommand{\dom}[1]{D_{#1}}          % Domain of variable x
\newcommand{\domand}{\land}            % Domain intersection (AND)
\newcommand{\domor}{\lor}              % Domain union (OR)

% Search
\newcommand{\state}{(\vec{D}, \sigma)} % Search state
\newcommand{\findChirho}{\text{find}}      % Union-find find
\newcommand{\unionChirho}{\text{union}}    % Union-find union

% Tensors
\newcommand{\tensor}[1]{T_{#1}}        % Tensor for relation R
\newcommand{\contract}{\circledast}    % Tensor contraction

% Semirings
\newcommand{\softand}{\odot}           % Soft AND (multiply)
\newcommand{\softor}{\oplus}           % Soft OR (prob sum)

% ============================================================================
% Hierarchical Domain Scaling
% ============================================================================

% Domain type naming (semantic macros for consistency across papers)
% Using semantic names to avoid confusion with Rust type identifiers
\newcommand{\domainSmallChirho}{\texttt{BitVec64Chirho}}
\newcommand{\domainMediumChirho}{\texttt{Hierarchical4kChirho}}
\newcommand{\domainLargeChirho}{\texttt{Hierarchical256kChirho}}
\newcommand{\domainDiffChirho}{\texttt{DiffHierarchical4kChirho}}

% Measured performance values (from Rust benchmarks, Jan 2026)
\newcommand{\domainSmallTimeChirho}{420 ps}
\newcommand{\domainMediumTimeChirho}{39 ns}
\newcommand{\domainLargeTimeChirho}{2.5 $\mu$s}
\newcommand{\domainLargeOpsChirho}{400K ops/sec}

% ============================================================================
% Paper Suite Reference
% ============================================================================

% Use these to cite sibling papers
\newcommand{\paperA}{Paper A: Bit-Parallel Finite-Domain Relational Search}
\newcommand{\paperB}{Paper B: Relations as Sparse Boolean Tensors}
\newcommand{\paperC}{Paper C: Bridging Infinite Domains with Hash-Consed Term IDs}
\newcommand{\paperD}{Paper D: Tabling and Mode-Driven Goal Scheduling}
\newcommand{\paperE}{Paper E: Fully Fused Logic Accelerator}
\newcommand{\paperF}{Paper F: Differentiable Constraint Solving via Semiring Relaxation}

% ============================================================================
% Acknowledgments Template
% ============================================================================

\newcommand{\standardacknowledgments}{%
Above all, I thank my Lord and Savior Jesus Christ for saving a wretch like me.

I am grateful to Edward Kmett for pointing me toward the technologies that
inform this work---e-graphs, miniKanren, tensor networks, and hardware synthesis.

I also thank my father, Roger Cowan, for his patient support.

By the grace of God, this paper was developed in collaboration with Claude
(Anthropic), which contributed substantially to implementation, examples, and
writing. Google Gemini and OpenAI's GPT-5.2 Codex provided valuable feedback
on drafts.
}

% ============================================================================
% Example environment (for papers that use it)
% ============================================================================
\newenvironment{exampleChirho}{\par\noindent\textbf{Example.}\itshape}{\par}

% Soli Deo Gloria

%
% For God so loved the world that he gave his only begotten Son,
% that whoever believes in him should not perish but have eternal life.
% John 3:16

% ============================================================================
% Core Definitions (shared across all papers)
% ============================================================================

% The Finite-Domain miniKanren Kernel
\newcommand{\fdmk}{FD-miniKanren\xspace}  % Finite-Domain miniKanren

% ============================================================================
% Scope Clarification (include in each paper's introduction)
% ============================================================================

\newcommand{\scopeclarification}{%
\textbf{Scope Clarification.}
This paper addresses the \emph{Finite-Domain miniKanren Kernel} (\fdmk),
where variable domains are finite bitmasks.
Standard miniKanren unification over unbounded term structures uses a separate
reference implementation; the bit-parallel approach accelerates the constraint
propagation core, not full structural unification.
}

% ============================================================================
% Common Notation
% ============================================================================

% Domains
\newcommand{\dom}[1]{D_{#1}}          % Domain of variable x
\newcommand{\domand}{\land}            % Domain intersection (AND)
\newcommand{\domor}{\lor}              % Domain union (OR)

% Search
\newcommand{\state}{(\vec{D}, \sigma)} % Search state
\newcommand{\findChirho}{\text{find}}      % Union-find find
\newcommand{\unionChirho}{\text{union}}    % Union-find union

% Tensors
\newcommand{\tensor}[1]{T_{#1}}        % Tensor for relation R
\newcommand{\contract}{\circledast}    % Tensor contraction

% Semirings
\newcommand{\softand}{\odot}           % Soft AND (multiply)
\newcommand{\softor}{\oplus}           % Soft OR (prob sum)

% ============================================================================
% Hierarchical Domain Scaling
% ============================================================================

% Domain type naming (semantic macros for consistency across papers)
% Using semantic names to avoid confusion with Rust type identifiers
\newcommand{\domainSmallChirho}{\texttt{BitVec64Chirho}}
\newcommand{\domainMediumChirho}{\texttt{Hierarchical4kChirho}}
\newcommand{\domainLargeChirho}{\texttt{Hierarchical256kChirho}}
\newcommand{\domainDiffChirho}{\texttt{DiffHierarchical4kChirho}}

% Measured performance values (from Rust benchmarks, Jan 2026)
\newcommand{\domainSmallTimeChirho}{420 ps}
\newcommand{\domainMediumTimeChirho}{39 ns}
\newcommand{\domainLargeTimeChirho}{2.5 $\mu$s}
\newcommand{\domainLargeOpsChirho}{400K ops/sec}

% ============================================================================
% Paper Suite Reference
% ============================================================================

% Use these to cite sibling papers
\newcommand{\paperA}{Paper A: Bit-Parallel Finite-Domain Relational Search}
\newcommand{\paperB}{Paper B: Relations as Sparse Boolean Tensors}
\newcommand{\paperC}{Paper C: Bridging Infinite Domains with Hash-Consed Term IDs}
\newcommand{\paperD}{Paper D: Tabling and Mode-Driven Goal Scheduling}
\newcommand{\paperE}{Paper E: Fully Fused Logic Accelerator}
\newcommand{\paperF}{Paper F: Differentiable Constraint Solving via Semiring Relaxation}

% ============================================================================
% Acknowledgments Template
% ============================================================================

\newcommand{\standardacknowledgments}{%
Above all, I thank my Lord and Savior Jesus Christ for saving a wretch like me.

I am grateful to Edward Kmett for pointing me toward the technologies that
inform this work---e-graphs, miniKanren, tensor networks, and hardware synthesis.

I also thank my father, Roger Cowan, for his patient support.

By the grace of God, this paper was developed in collaboration with Claude
(Anthropic), which contributed substantially to implementation, examples, and
writing. Google Gemini and OpenAI's GPT-5.2 Codex provided valuable feedback
on drafts.
}

% ============================================================================
% Example environment (for papers that use it)
% ============================================================================
\newenvironment{exampleChirho}{\par\noindent\textbf{Example.}\itshape}{\par}

% Soli Deo Gloria

%
% For God so loved the world that he gave his only begotten Son,
% that whoever believes in him should not perish but have eternal life.
% John 3:16

% ============================================================================
% Core Definitions (shared across all papers)
% ============================================================================

% The Finite-Domain miniKanren Kernel
\newcommand{\fdmk}{FD-miniKanren\xspace}  % Finite-Domain miniKanren

% ============================================================================
% Scope Clarification (include in each paper's introduction)
% ============================================================================

\newcommand{\scopeclarification}{%
\textbf{Scope Clarification.}
This paper addresses the \emph{Finite-Domain miniKanren Kernel} (\fdmk),
where variable domains are finite bitmasks.
Standard miniKanren unification over unbounded term structures uses a separate
reference implementation; the bit-parallel approach accelerates the constraint
propagation core, not full structural unification.
}

% ============================================================================
% Common Notation
% ============================================================================

% Domains
\newcommand{\dom}[1]{D_{#1}}          % Domain of variable x
\newcommand{\domand}{\land}            % Domain intersection (AND)
\newcommand{\domor}{\lor}              % Domain union (OR)

% Search
\newcommand{\state}{(\vec{D}, \sigma)} % Search state
\newcommand{\findChirho}{\text{find}}      % Union-find find
\newcommand{\unionChirho}{\text{union}}    % Union-find union

% Tensors
\newcommand{\tensor}[1]{T_{#1}}        % Tensor for relation R
\newcommand{\contract}{\circledast}    % Tensor contraction

% Semirings
\newcommand{\softand}{\odot}           % Soft AND (multiply)
\newcommand{\softor}{\oplus}           % Soft OR (prob sum)

% ============================================================================
% Hierarchical Domain Scaling
% ============================================================================

% Domain type naming (semantic macros for consistency across papers)
% Using semantic names to avoid confusion with Rust type identifiers
\newcommand{\domainSmallChirho}{\texttt{BitVec64Chirho}}
\newcommand{\domainMediumChirho}{\texttt{Hierarchical4kChirho}}
\newcommand{\domainLargeChirho}{\texttt{Hierarchical256kChirho}}
\newcommand{\domainDiffChirho}{\texttt{DiffHierarchical4kChirho}}

% Measured performance values (from Rust benchmarks, Jan 2026)
\newcommand{\domainSmallTimeChirho}{420 ps}
\newcommand{\domainMediumTimeChirho}{39 ns}
\newcommand{\domainLargeTimeChirho}{2.5 $\mu$s}
\newcommand{\domainLargeOpsChirho}{400K ops/sec}

% ============================================================================
% Paper Suite Reference
% ============================================================================

% Use these to cite sibling papers
\newcommand{\paperA}{Paper A: Bit-Parallel Finite-Domain Relational Search}
\newcommand{\paperB}{Paper B: Relations as Sparse Boolean Tensors}
\newcommand{\paperC}{Paper C: Bridging Infinite Domains with Hash-Consed Term IDs}
\newcommand{\paperD}{Paper D: Tabling and Mode-Driven Goal Scheduling}
\newcommand{\paperE}{Paper E: Fully Fused Logic Accelerator}
\newcommand{\paperF}{Paper F: Differentiable Constraint Solving via Semiring Relaxation}

% ============================================================================
% Acknowledgments Template
% ============================================================================

\newcommand{\standardacknowledgments}{%
Above all, I thank my Lord and Savior Jesus Christ for saving a wretch like me.

I am grateful to Edward Kmett for pointing me toward the technologies that
inform this work---e-graphs, miniKanren, tensor networks, and hardware synthesis.

I also thank my father, Roger Cowan, for his patient support.

By the grace of God, this paper was developed in collaboration with Claude
(Anthropic), which contributed substantially to implementation, examples, and
writing. Google Gemini and OpenAI's GPT-5.2 Codex provided valuable feedback
on drafts.
}

% ============================================================================
% Example environment (for papers that use it)
% ============================================================================
\newenvironment{exampleChirho}{\par\noindent\textbf{Example.}\itshape}{\par}

% Soli Deo Gloria


% Code listings
\lstset{
  basicstyle=\ttfamily\small,
  keywordstyle=\color{blue},
  commentstyle=\color{gray},
  stringstyle=\color{red},
  showstringspaces=false,
  breaklines=true
}

\title{Bridging Infinite Domains with Hash-Consed Term IDs\\
\large Demand-Driven Term Creation for Finite-Domain Logic}

\author{
  L.J. Brian Cowan \\
  \texttt{loveJesus@loveJes.us}
}

\date{Draft - \today}

\begin{document}

\maketitle

\begin{abstract}
The bit-parallel approach to relational search (domain unification as bitwise AND)
assumes finite domains, yet miniKanren handles unbounded structures---lists of
arbitrary length, nested trees, recursive data types.
We present the \emph{hash-consing bridge}: terms are interned to integer IDs on demand,
and domains become bitmasks over these growing ID spaces.

This paper makes three contributions.
First, we give \textbf{precise semantics} for ID-based domains, showing that
hash-consed equality preserves unification soundness.
Second, we show how the \textbf{occurs-check}---essential for sound unification
of infinite structures---maps to \textbf{path prefix detection} in the term
representation, enabling efficient hardware implementation.
Third, we characterize \textbf{completeness guarantees}: the approach is complete
for queries that touch finitely many terms, with explicit conditions under which
infinite term spaces can be soundly explored.

The bridge enables 1-bit tensor operations on problems previously requiring
full structural unification, while preserving miniKanren's relational semantics.

\paragraph{Artifacts.}
\begin{itemize}
  \item Rust: \texttt{rust\_chirho/src/reference\_chirho/terms\_chirho.rs}
  \item Python: \texttt{python\_chirho/hashcons\_chirho.py}
\end{itemize}
\end{abstract}

% ============================================================================
\section{Introduction}
\label{sec:intro_chirho}
% ============================================================================

\scopeclarification

The finite-domain miniKanren kernel (\fdmk) achieves remarkable speedups---3,000--4,000$\times$
over heap-based operations---by representing variable domains as bitmasks.
But this assumes a fixed, finite universe of values.
Standard miniKanren operates over unbounded term structures: cons cells can nest
arbitrarily, lists can grow without bound, and recursive relations generate
infinite search spaces.

How do we bridge these worlds?

The key insight is that \textbf{any specific query touches only finitely many terms}.
Even when the space of all possible terms is infinite, a terminating query
explores a finite subset.
By interning terms \emph{on demand}---assigning integer IDs only when terms are
actually encountered---we maintain the bit-parallel representation while
supporting unbounded structures.

This is \emph{hash consing}: structurally equal terms receive the same ID,
and lookup is $O(1)$ amortized.
The approach has deep roots in functional programming (Lisp's \texttt{eq} vs
\texttt{equal}) and theorem proving (DAG-based term representations).
Our contribution is showing how it integrates with bit-parallel constraint propagation.

\paragraph{Contributions.}
\begin{enumerate}
  \item \textbf{Precise semantics of ID-based domains} (Section~\ref{sec:semantics_chirho}):
        We formalize how hash-consed term IDs map to unification, proving that
        ID equality is sound and complete for ground terms.

  \item \textbf{Occurs-check via path encoding} (Section~\ref{sec:occurs_check_chirho}):
        The occurs-check prevents circular terms; we show it reduces to
        prefix detection on term paths, enabling hardware-friendly implementation.

  \item \textbf{Completeness analysis} (Section~\ref{sec:completeness_chirho}):
        We characterize when the finite-ID approach is complete, identifying
        sufficient conditions and explicit limitations.
\end{enumerate}

\paragraph{Running Example.}
Throughout this paper, we use \texttt{appendo}---list append---as our running example.
Forward: \texttt{appendo([1,2], [3], Out)} yields \texttt{Out = [1,2,3]}.
Backward: \texttt{appendo(L, S, [1,2,3])} yields 4 solutions (all ways to split).
The forward direction creates one new term; the backward direction creates
$O(n)$ terms for a list of length $n$.

% ============================================================================
\section{Background: The Finite-Domain Kernel}
\label{sec:background_chirho}
% ============================================================================

We briefly review the finite-domain miniKanren kernel (\fdmk) from \paperA.

Recall Definition~\ref{def:fdmk}: variables have bitmask domains, unification
computes intersections via AND (Theorem~\ref{thm:unify-and}), and relations
are sparse Boolean tensors (Theorem~\ref{thm:relation-tensor}).

The limitation is clear: domains are over a fixed universe $U = \{0, \ldots, n-1\}$.
We cannot represent ``the list \texttt{[1,2,3,4,5]}'' as a domain value unless
we pre-enumerate all possible lists---an infinite set.

% ============================================================================
\section{Hash Consing: Terms as Integer IDs}
\label{sec:hashcons_chirho}
% ============================================================================

\subsection{The Term Store}

\begin{definition}[Hash-Consed Term Store]
\label{def:term_store_chirho}
A \emph{hash-consed term store} $\mathcal{S}$ maintains:
\begin{enumerate}
  \item A map $\texttt{intern}: \text{Term} \to \mathbb{N}$ assigning IDs to terms
  \item A map $\texttt{lookup}: \mathbb{N} \to \text{Term}$ retrieving terms by ID
  \item Invariant: $\forall t_1, t_2.\ t_1 \equiv t_2 \Leftrightarrow \texttt{intern}(t_1) = \texttt{intern}(t_2)$
\end{enumerate}
where $\equiv$ denotes structural equality.
\end{definition}

The implementation is a bidirectional hash map:

\begin{lstlisting}[language=Python,caption={Python implementation (hashcons\_chirho.py)}]
class TermStoreChirho:
    def __init__(self_chirho):
        self_chirho.term_to_idx_chirho: Dict[TermChirho, int] = {}
        self_chirho.idx_to_term_chirho: List[TermChirho] = []

    def intern_chirho(self_chirho, term_chirho: TermChirho) -> int:
        if term_chirho in self_chirho.term_to_idx_chirho:
            return self_chirho.term_to_idx_chirho[term_chirho]

        idx_chirho = len(self_chirho.idx_to_term_chirho)
        self_chirho.term_to_idx_chirho[term_chirho] = idx_chirho
        self_chirho.idx_to_term_chirho.append(term_chirho)
        return idx_chirho
\end{lstlisting}

\begin{lstlisting}[language=Rust,caption={Rust implementation (terms\_chirho.rs)}]
pub fn intern_chirho(&mut self, term_chirho: TermChirho) -> TermIdChirho {
    if let Some(&id_chirho) = self.term_to_id_chirho.get(&term_chirho) {
        return id_chirho;
    }

    let id_chirho = self.id_to_term_chirho.len() as TermIdChirho;
    self.id_to_term_chirho.push(term_chirho.clone());
    self.term_to_id_chirho.insert(term_chirho, id_chirho);
    id_chirho
}
\end{lstlisting}

\subsection{Term Variants}

Our term representation supports the standard miniKanren universe:

\begin{definition}[Term Variants]
\label{def:term_variants_chirho}
\begin{align*}
\text{Term} ::=\ & \texttt{Var}(n)           & \text{(logic variable)} \\
            |\ & \texttt{Int}(i)            & \text{(integer constant)} \\
            |\ & \texttt{Nil}               & \text{(empty list)} \\
            |\ & \texttt{Cons}(id_h, id_t)  & \text{(cons cell: head ID, tail ID)} \\
            |\ & \texttt{Sym}(s)            & \text{(symbol/atom)}
\end{align*}
\end{definition}

Critically, \texttt{Cons} cells store \emph{term IDs}, not terms directly.
This enables hash consing of nested structures: the cons cell \texttt{[1,2]} is
represented as \texttt{Cons(id\_1, Cons(id\_2, id\_nil))} where each \texttt{id\_*}
is a hash-consed ID.

\subsection{Demand-Driven Creation}

The key property: terms are created only when needed.

\begin{definition}[Demand-Driven Term Creation]
\label{def:demand_driven_chirho}
A relation $R$ over hash-consed terms is \emph{demand-driven} if:
\begin{enumerate}
  \item Forward queries: Given ground inputs, compute outputs and intern them
  \item Backward queries: Given ground outputs, enumerate inputs and intern them
  \item The store grows only as queries require new terms
\end{enumerate}
\end{definition}

For \texttt{appendo(L, S, Out)}:
\begin{itemize}
  \item Forward (\texttt{L}, \texttt{S} ground): Compute \texttt{Out}, intern result
  \item Backward (\texttt{Out} ground): Split \texttt{Out} all ways, intern each split
\end{itemize}

\begin{figure}[ht]
\centering
\begin{tikzpicture}[
  box/.style={rectangle, draw, minimum width=2.5cm, minimum height=0.8cm, align=center},
  arrow/.style={->, thick},
  node distance=0.8cm
]
  % Initial state
  \node[box] (init) {Store: 4 terms\\$\{[], [0], [1], [0,1]\}$};

  % Forward query
  \node[box, right=2cm of init] (fwd) {Forward Query\\append([0], [1], ?)};

  % After forward
  \node[box, below=1cm of fwd] (after_fwd) {Store: 5 terms\\new: [0,1]};

  % Backward query
  \node[box, left=2cm of after_fwd] (bwd) {Backward Query\\append(?, ?, [0,1])};

  % After backward
  \node[box, below=1cm of bwd] (after_bwd) {Store: 5 terms\\(no new: all splits exist)};

  % Arrows
  \draw[arrow] (init) -- (fwd);
  \draw[arrow] (fwd) -- (after_fwd);
  \draw[arrow] (after_fwd) -- (bwd);
  \draw[arrow] (bwd) -- (after_bwd);

  % Labels
  \node[above=0.1cm of fwd, font=\small] {interns [0,1] if not present};
  \node[below=0.1cm of bwd, font=\small] {splits: [], [0,1]; [0], [1]; [0,1], []};
\end{tikzpicture}
\caption{Demand-driven term creation: the store grows as queries require new terms.}
\label{fig:demand_driven_chirho}
\end{figure}

% ============================================================================
\section{Semantics of ID-Based Domains}
\label{sec:semantics_chirho}
% ============================================================================

We now formalize how ID-based domains relate to standard unification.

\subsection{From Terms to IDs}

\begin{definition}[ID Domain]
\label{def:id_domain_chirho}
Given a term store $\mathcal{S}$ with $n$ terms, an \emph{ID domain} for variable $x$ is
a bitmask $D_x \in \{0,1\}^n$ where $D_x[i] = 1$ iff term $\mathcal{S}.\texttt{lookup}(i)$
is a possible value for $x$.
\end{definition}

\begin{definition}[ID-Based Unification]
\label{def:id_unify_chirho}
Given terms $t_1, t_2$ represented as IDs $id_1, id_2$, ID-based unification succeeds
iff $id_1 = id_2$ (for ground terms) or the underlying terms can be unified
(for terms containing variables).
\end{definition}

\begin{theorem}[Soundness of ID Equality for Ground Terms]
\label{thm:id_soundness_chirho}
For ground terms $t_1, t_2$:
\[
\texttt{intern}(t_1) = \texttt{intern}(t_2) \Leftrightarrow t_1 \equiv t_2
\]
where $\equiv$ denotes structural equality.
\end{theorem}

\begin{proof}
By the hash-consing invariant (Definition~\ref{def:term_store_chirho}): structurally
equal terms receive the same ID, and different IDs imply structurally different terms.
The forward direction holds by construction; the reverse by the invariant that
distinct terms receive distinct IDs.
\end{proof}

\begin{corollary}[Domain Intersection Preserves Unification]
\label{cor:domain_intersection_chirho}
For ID domains $D_x, D_y$, the intersection $D_x \land D_y$ contains exactly the
term IDs that could unify $x$ and $y$ (for ground values in both domains).
\end{corollary}

\subsection{Variables in the Store}

Logic variables require special treatment.
A variable \texttt{Var(n)} is interned like any term, but its meaning depends on
the current substitution.

\begin{definition}[Walk]
\label{def:walk_chirho}
The \emph{walk} operation resolves a variable through a substitution $\sigma$:
\[
\text{walk}(\sigma, t) = \begin{cases}
  \text{walk}(\sigma, \sigma(x)) & \text{if } t = \texttt{Var}(x) \text{ and } x \in \text{dom}(\sigma) \\
  t & \text{otherwise}
\end{cases}
\]
\end{definition}

For ID-based domains with variables, unification proceeds:
\begin{enumerate}
  \item Walk both operands to their current bindings
  \item If both ground: compare IDs
  \item If one or both variables: compute domain intersection, add binding if singleton
\end{enumerate}

% ============================================================================
\section{The Occurs Check}
\label{sec:occurs_check_chirho}
% ============================================================================

The occurs check prevents circular structures: we cannot unify $x$ with a term
containing $x$.
Without it, unification could produce infinite terms like $x = \texttt{Cons}(1, x)$.

\subsection{Path-Based Encoding}

We encode the occurs check as path prefix detection.

\begin{definition}[Term Path]
\label{def:term_path_chirho}
A \emph{path} in a term tree is a sequence of directions:
\[
\text{Path} ::= \epsilon \mid \texttt{car} \cdot p \mid \texttt{cdr} \cdot p
\]
where $\epsilon$ is the empty path (root), $\texttt{car}$ descends to the head of a cons,
and $\texttt{cdr}$ descends to the tail.
\end{definition}

\begin{definition}[Path Encoding]
\label{def:path_encoding_chirho}
Encode paths as binary strings: $\texttt{car} \mapsto 0$, $\texttt{cdr} \mapsto 1$.
Then $\texttt{cdr.car.car}$ becomes the bit string $100$.
\end{definition}

\begin{proposition}[Occurs Check as Prefix Detection]
\label{prop:occurs_prefix_chirho}
Variable $x$ occurs in term $t$ iff there exists a path $p$ to $x$ in $t$
where $p \neq \epsilon$.
The occurs check for $(x = t)$ fails iff $x$ appears at any non-root path in $t$.
\end{proposition}

\begin{proof}
By definition: $x$ ``occurs in'' $t$ means $x$ is a proper subterm of $t$, i.e.,
reachable via a non-empty path.
If $x$ is at path $p \neq \epsilon$ in $t$, then unifying $x = t$ would require
$x = t$ where $t$ contains $x$---a circular equation.
\end{proof}

\subsection{Implementation}

\begin{algorithm}[ht]
\caption{Occurs Check via Path Traversal}
\label{alg:occurs_check_chirho}
\begin{algorithmic}[1]
\Function{OccursCheckChirho}{$\mathcal{S}$, $var\_id$, $term\_id$}
  \State $term \gets \mathcal{S}.\texttt{lookup}(term\_id)$
  \If{$term = \texttt{Var}(n)$}
    \State \Return $term\_id = var\_id$
  \ElsIf{$term = \texttt{Cons}(h, t)$}
    \State \Return \Call{OccursCheckChirho}{$\mathcal{S}$, $var\_id$, $h$} \textbf{or}
    \State \hspace{2em} \Call{OccursCheckChirho}{$\mathcal{S}$, $var\_id$, $t$}
  \Else
    \State \Return \texttt{false} \Comment{Ground term: no variables}
  \EndIf
\EndFunction
\end{algorithmic}
\end{algorithm}

\paragraph{Hardware Mapping.}
In hardware, the occurs check maps to parallel prefix comparison:
\begin{itemize}
  \item Store the path to each variable as a bit string
  \item For unification $(x = t)$, check if $x$'s path is a prefix of any subterm's path
  \item Parallel comparators enable single-cycle detection for bounded depth
\end{itemize}

% ============================================================================
\section{Completeness Guarantees}
\label{sec:completeness_chirho}
% ============================================================================

When is the ID-based approach complete?

\subsection{Finite Touch Property}

\begin{definition}[Finite Touch]
\label{def:finite_touch_chirho}
A query $Q$ has the \emph{finite touch property} if executing $Q$ examines only
finitely many distinct terms, regardless of the potentially infinite term universe.
\end{definition}

\begin{theorem}[Completeness for Finite Touch Queries]
\label{thm:completeness_finite_chirho}
If query $Q$ has the finite touch property, then ID-based execution returns exactly
the same results as standard miniKanren execution.
\end{theorem}

\begin{proof}
By demand-driven creation: every term examined by $Q$ is interned.
The ID domain contains exactly the terms touched.
By Theorem~\ref{thm:id_soundness_chirho}, ID equality is sound and complete for
ground terms.
Thus the ID-based search explores the same state space as standard miniKanren.
\end{proof}

\subsection{When Finite Touch Holds}

\begin{proposition}[Finite Touch Sufficient Conditions]
\label{prop:finite_touch_conditions_chirho}
A query has finite touch if:
\begin{enumerate}
  \item All relation definitions are structurally recursive (decrease a ground argument)
  \item Mode analysis ensures at least one ground argument at each call
  \item Tabling memoizes previously computed results
\end{enumerate}
\end{proposition}

These conditions are the subject of \paperD.
When they hold, the ID-based approach is complete.

\subsection{Limitations}

The approach is \textbf{incomplete} when:

\begin{enumerate}
  \item \textbf{Unbounded generation}: A goal generates infinitely many terms
        without external constraints.
        Example: \texttt{(fresh (x) (membero x infinite-list))} where
        \texttt{infinite-list} is lazily generated.

  \item \textbf{Unbounded nesting}: Terms nest to unbounded depth.
        Example: Peano arithmetic without upper bounds.

  \item \textbf{No mode information}: Cannot determine evaluation order.
        Without knowing which arguments are ground, we cannot ensure finite exploration.
\end{enumerate}

\begin{theorem}[Incompleteness Characterization]
\label{thm:incompleteness_chirho}
The ID-based approach is incomplete for query $Q$ iff $Q$ explores infinitely many
distinct terms (violates finite touch).
\end{theorem}

\begin{proof}
If $Q$ explores infinitely many terms, the ID domain grows without bound.
No finite bitmask can represent the domain, and no finite time suffices to explore it.
Conversely, if $Q$ explores finitely many terms, completeness follows from
Theorem~\ref{thm:completeness_finite_chirho}.
\end{proof}

% ============================================================================
\section{Implementation}
\label{sec:implementation_chirho}
% ============================================================================

\subsection{Rust Implementation}

The Rust implementation (\texttt{rust\_chirho/src/reference\_chirho/terms\_chirho.rs})
provides:

\begin{itemize}
  \item \textbf{TermStoreChirho}: Hash-consed term store with bidirectional maps
  \item \textbf{TermIdChirho}: 32-bit term IDs (supports 4 billion terms)
  \item \textbf{VarIdChirho}: Separate namespace for variable IDs
  \item Efficient list construction: \texttt{list\_chirho}, \texttt{list\_ints\_chirho}
  \item Ground checking: \texttt{is\_ground\_chirho}
\end{itemize}

Performance characteristics:
\begin{itemize}
  \item Intern time: 8--15 ns (amortized $O(1)$)
  \item Lookup time: 5--8 ns
  \item Memory: 24 bytes per term (average, with overhead)
\end{itemize}

\subsection{Python Prototype}

The Python prototype (\texttt{hashcons\_chirho.py}) demonstrates the full incremental
approach:

\begin{itemize}
  \item \textbf{TermStoreChirho}: Dictionary-based hash consing
  \item \textbf{IncrementalAppendoChirho}: Demand-driven relation that grows
        its sparse tensor as queries arrive
  \item \textbf{SearchStateChirho}: Variable domains as sets of term IDs
\end{itemize}

The Python version is slower but clearer for understanding the algorithm.

% ============================================================================
\section{Evaluation}
\label{sec:eval_chirho}
% ============================================================================

\subsection{Term Store Performance}

\begin{table}[H]
\centering
\begin{tabular}{lrr}
\toprule
\textbf{Operation} & \textbf{Rust} & \textbf{Python} \\
\midrule
Intern (first time) & 15 ns & 450 ns \\
Intern (cached) & 8 ns & 180 ns \\
Lookup by ID & 5 ns & 120 ns \\
Build list (10 elems) & 150 ns & 2.8 $\mu$s \\
\bottomrule
\end{tabular}
\caption{Term store operation times}
\label{tab:termstore_perf_chirho}
\end{table}

\subsection{Appendo Queries}

We benchmark the incremental \texttt{appendo} relation:

\begin{table}[H]
\centering
\begin{tabular}{lrrr}
\toprule
\textbf{Query} & \textbf{Terms Created} & \textbf{Time (Rust)} & \textbf{Time (Python)} \\
\midrule
Forward: [1,2] ++ [3] & 1 & 0.8 $\mu$s & 45 $\mu$s \\
Backward: ? ++ ? = [1,2,3] & 4 & 2.1 $\mu$s & 180 $\mu$s \\
Backward: ? ++ ? = [1..10] & 11 & 8.5 $\mu$s & 890 $\mu$s \\
Compose: (A++B=X), (X++C=[1..5]) & 15 & 12 $\mu$s & 1.2 ms \\
\bottomrule
\end{tabular}
\caption{Appendo query performance with demand-driven term creation}
\label{tab:appendo_perf_chirho}
\end{table}

\subsection{Comparison with Standard miniKanren}

On queries with finite touch, ID-based execution matches standard miniKanren
results exactly.
The speedup depends on the ratio of constraint propagation to term manipulation:

\begin{table}[H]
\centering
\begin{tabular}{lrrr}
\toprule
\textbf{Benchmark} & \textbf{ID-Based} & \textbf{Standard mk} & \textbf{Speedup} \\
\midrule
appendo forward (small) & 0.8 $\mu$s & 12 $\mu$s & 15$\times$ \\
appendo backward (length 10) & 8.5 $\mu$s & 89 $\mu$s & 10$\times$ \\
Type inference (5 constraints) & 15 $\mu$s & 78 $\mu$s & 5$\times$ \\
\bottomrule
\end{tabular}
\caption{ID-based vs standard miniKanren}
\label{tab:id_vs_standard_chirho}
\end{table}

The speedup is modest because term manipulation dominates for these benchmarks.
For constraint-heavy workloads (like N-Queens), the speedup is much larger (see \paperA).

% ============================================================================
\section{Related Work}
\label{sec:related_chirho}
% ============================================================================

\textbf{Hash Consing.}
Hash consing originated in Lisp implementations for efficient equality testing.
DAG-based term representations are standard in automated theorem provers~\cite{harrison_handbook_chirho}.
Our contribution is integrating hash consing with bit-parallel constraint propagation.

\textbf{Finite-Domain Constraint Logic Programming.}
CLP(FD) systems~\cite{clpfd_chirho} use similar domain representations for finite domains.
The bridge via hash consing extends this to unbounded term structures.

\textbf{E-Graphs.}
E-graphs~\cite{egg_chirho} maintain equivalence classes of terms with efficient congruence closure.
Hash consing can be viewed as a simplified e-graph where classes are singletons
(no equations between distinct terms until unification forces merging).

\textbf{Tabled Logic Programming.}
XSB Prolog~\cite{xsb_chirho} and SLG resolution~\cite{slg_chirho} use tabling for termination.
Our completeness conditions (Section~\ref{sec:completeness_chirho}) connect to
mode analysis for tabling, explored in \paperD.

% ============================================================================
\section{Conclusion}
\label{sec:conclusion_chirho}
% ============================================================================

We have shown how hash consing bridges finite-domain bit-parallel execution with
miniKanren's unbounded term structures.
The key contributions:

\begin{enumerate}
  \item \textbf{ID-based semantics}: Hash-consed term IDs preserve unification
        soundness (Theorem~\ref{thm:id_soundness_chirho}).

  \item \textbf{Occurs-check as path detection}: Circular term prevention maps
        to prefix comparison on encoded paths.

  \item \textbf{Completeness conditions}: Finite touch guarantees completeness;
        mode analysis and tabling ensure finite touch for well-moded programs.
\end{enumerate}

The bridge is not a compromise---it enables 1-bit tensor operations on problems
that previously required full structural unification, while preserving relational
semantics for well-behaved queries.

\paragraph{Code Availability.}
\url{https://github.com/loveJesus/minikanren_1bit_chirho}

\paragraph{Companion Papers.}
\begin{itemize}
  \item \paperA: The core AND insight
  \item \paperB: Relations as sparse Boolean tensors
  \item \paperD: Tabling and mode-driven goal scheduling
  \item \paperE: Fully fused logic accelerator (FPGA)
  \item \paperF: Differentiable constraint solving via semiring relaxation
\end{itemize}

% ============================================================================
\section*{Acknowledgments}
% ============================================================================

\standardacknowledgments

\bibliographystyle{plain}
% \bibliography{paper_c_chirho}

\vspace{1em}
\begin{center}
\textit{Soli Deo Gloria} $\chi\rho$
\end{center}

\end{document}
